\documentclass[11pt,a4paper]{article}
\usepackage[margin=0.9in]{geometry}
\usepackage{amsmath,amssymb,amsthm}
\usepackage{enumitem}
\usepackage{array}
\usepackage{longtable}
\usepackage{booktabs}
\usepackage{xcolor}
\usepackage{titlesec}
\usepackage{fancyhdr}
\setlength{\headheight}{14pt}
\usepackage{hyperref}

\hypersetup{colorlinks=true, linkcolor=blue!60!black, urlcolor=blue!60!black}

\titleformat{\section}{\Large\bfseries\color{black!80}}{\thesection.}{0.6em}{}
\titleformat{\subsection}{\large\bfseries\color{black!70}}{\thesubsection}{0.6em}{}

\pagestyle{fancy}
\fancyhf{}
\fancyhead[L]{\textbf{CAT-Hardened QA Simulation Batch}}
\fancyhead[R]{\thepage}

\newcommand{\qhead}[4]{%
  \vspace{0.7em}\noindent\textbf{Q#1.}\hfill\textit{Topic: #2 \textbar{} Subtopic: #3 \textbar{} Difficulty: #4}\\[-0.2em]
}

\title{\textbf{CAT-Hardened Quantitative Aptitude Simulation}\\[0.3em]
\large Practice Batch -- 15 Questions}
\author{Exam-Grade Practice Set}
\date{}

\begin{document}
\maketitle

%================================================================
\section*{Instructions}

\begin{enumerate}[leftmargin=*,itemsep=0.2em]
  \item This batch contains \textbf{15 questions}, distributed across Arithmetic (5), Algebra (4), Number System (3), and Geometry/Mensuration (3).
  \item Question types: \textbf{9 MCQ} and \textbf{6 TITA} (Type In The Answer).
  \item For MCQ questions, exactly one option is correct.
  \item For TITA questions, enter the exact numerical value.
  \item Calculators are not permitted.
  \item Recommended time: \textbf{40--45 minutes}.
  \item In an actual CAT, negative marking applies on MCQs (1/3 of marks per wrong answer).
\end{enumerate}

\vspace{0.4em}\hrule\vspace{0.6em}

%================================================================
\section*{Question Paper}

%--------------------------- ARITHMETIC ---------------------------
\subsection*{Arithmetic}

\qhead{1}{Arithmetic}{Time, Speed and Distance}{Hard \textbar{} MCQ}
Two cyclists $P$ and $Q$ start simultaneously from towns $A$ and $B$ respectively and ride towards each other on a straight road at constant speeds. They first meet at a point $7.5$ km from $A$. After meeting, each continues to the other town, turns back immediately on reaching it, and they meet for the second time at a point $4$ km from $B$. The distance (in km) between $A$ and $B$ is:
\begin{enumerate}[label=(\Alph*),leftmargin=*]
  \item $17$
  \item $18.5$
  \item $19$
  \item $20$
\end{enumerate}

\qhead{2}{Arithmetic}{Mixtures and Replacement}{Very Hard \textbar{} TITA}
A vessel contains $80$ litres of pure milk. In each operation, exactly $x$ litres of the liquid in the vessel is removed and replaced by $x$ litres of water, where $x$ is the same positive integer in every operation. After exactly $4$ operations, the volume of milk remaining in the vessel is found to be strictly greater than $25$ litres and strictly less than $30$ litres. Find the largest possible value of $x$.

\qhead{3}{Arithmetic}{Time and Work}{Hard \textbar{} MCQ}
A and B together can finish a project in $20$ days. B alone takes $45$ days to finish the project. They start the project together, but A leaves after some days; B alone then takes a further $21.6$ days to complete the remaining work. The number of days A alone would take to complete the project is:
\begin{enumerate}[label=(\Alph*),leftmargin=*]
  \item $30$
  \item $36$
  \item $45$
  \item $60$
\end{enumerate}

\qhead{4}{Arithmetic}{Profit, Loss and Discount}{Medium-Hard \textbar{} TITA}
A trader marks his goods $40\%$ above cost price and offers a discount of $d\%$ on the marked price. He also uses a faulty balance that displays $1000$ g but actually delivers only $800$ g of the goods. If his overall profit is exactly $40\%$ of the actual cost price of the goods he physically delivers, find the value of $d$.

\qhead{5}{Arithmetic}{Averages and Ratios}{Hard \textbar{} MCQ}
In a class, the average mark of the boys is $68$ and that of the girls is $80$. The overall class average is $74$. A new group of students, consisting only of boys with average mark $62$, joins the class, and after they join, the overall class average becomes $71$. The ratio of the original number of boys in the class to the number of boys in the new group is:
\begin{enumerate}[label=(\Alph*),leftmargin=*]
  \item $1:1$
  \item $3:2$
  \item $2:1$
  \item $4:3$
\end{enumerate}

%--------------------------- ALGEBRA ---------------------------
\subsection*{Algebra}

\qhead{6}{Algebra}{Functional Equations}{Hard \textbar{} TITA}
Let $f:\mathbb{R}\to\mathbb{R}$ satisfy $f(x)+2\,f(1-x)=x^{2}+3x$ for every real $x$. Find the value of $3\,f(5)$.

\qhead{7}{Algebra}{Logarithms}{Hard \textbar{} MCQ}
The number of real values of $x$ satisfying
\[
\log_{2}(x^{2}-3x+2)\;+\;\log_{2}(x^{2}-7x+12)\;=\;3\;+\;\log_{2}\!\left(\frac{(x-1)(x-4)}{2}\right)
\]
is:
\begin{enumerate}[label=(\Alph*),leftmargin=*]
  \item $0$
  \item $1$
  \item $2$
  \item Infinitely many
\end{enumerate}

\qhead{8}{Algebra}{Inequalities with Constraints}{Very Hard \textbar{} MCQ}
Let $a, b, c$ be positive real numbers satisfying $a+b+c=12$ and $abc=64$. The maximum possible value of $a$ is:
\begin{enumerate}[label=(\Alph*),leftmargin=*]
  \item $4$
  \item $4+4\sqrt{2}$
  \item $8$
  \item $4+2\sqrt{3}$
\end{enumerate}

\qhead{9}{Algebra}{Quadratic Equations}{Medium-Hard \textbar{} MCQ}
Let $p$ and $q$ be the roots of the equation $x^{2}-(k+3)\,x+(2k+1)=0$. If $(p-q)^{2}=12$, then the sum of all possible values of $k$ equals:
\begin{enumerate}[label=(\Alph*),leftmargin=*]
  \item $-2$
  \item $2$
  \item $7$
  \item $-7$
\end{enumerate}

%--------------------------- NUMBER SYSTEM ---------------------------
\subsection*{Number System}

\qhead{10}{Number System}{Remainders / CRT}{Very Hard \textbar{} TITA}
Let $N$ be the smallest positive integer such that $N$ leaves remainder $3$ when divided by $7$, remainder $5$ when divided by $11$, and is divisible by $13$. If $M$ is the next positive integer (after $N$) satisfying all three conditions, find the value of $M-N$.

\qhead{11}{Number System}{Divisibility and Inclusion--Exclusion}{Hard \textbar{} MCQ}
The number of positive integers $n \le 600$ such that $n$ is divisible by exactly two of the three numbers $4$, $6$, and $9$ is:
\begin{enumerate}[label=(\Alph*),leftmargin=*]
  \item $17$
  \item $34$
  \item $51$
  \item $67$
\end{enumerate}

\qhead{12}{Number System}{Digits and Divisibility}{Hard \textbar{} TITA}
A three-digit positive integer $\overline{abc}$ (with $a \neq 0$) satisfies all of the following simultaneously:
\begin{enumerate}[label=(\roman*),leftmargin=*,itemsep=0pt]
  \item it is divisible by $11$;
  \item the sum of its digits is $14$;
  \item the number $\overline{cba}$ formed by reversing its digits (where $c\neq 0$) is strictly greater than $\overline{abc}$.
\end{enumerate}
Find the sum of all three-digit numbers $\overline{abc}$ satisfying these conditions.

%--------------------------- GEOMETRY ---------------------------
\subsection*{Geometry / Mensuration}

\qhead{13}{Geometry}{Triangles -- Cevians and Ratios}{Hard \textbar{} MCQ}
In triangle $ABC$, point $D$ lies on side $BC$ such that $BD:DC=2:3$. Point $E$ lies on segment $AD$ such that $AE:ED=3:2$. The line through $B$ and $E$, when extended, meets side $AC$ at point $F$. The ratio $AF:FC$ equals:
\begin{enumerate}[label=(\Alph*),leftmargin=*]
  \item $3:5$
  \item $3:4$
  \item $9:10$
  \item $2:3$
\end{enumerate}

\qhead{14}{Mensuration}{Circles and Tangents}{Hard \textbar{} TITA}
Two circles of radii $5$ and $3$ are externally tangent to each other. A line tangent externally to both circles touches the larger circle at $P$ and the smaller circle at $Q$. Find the length of $PQ$ (in the form $a\sqrt{b}$, enter the value of $a\cdot b$; for example, if $PQ=2\sqrt{15}$, enter $30$).

\qhead{15}{Geometry}{Coordinate Geometry and Optimisation}{Very Hard \textbar{} MCQ}
In the coordinate plane, $A=(0,0)$ and $B=(10,0)$. A point $P=(x,y)$ with $y>0$ is chosen such that the area of triangle $APB$ is exactly $20$ square units. Over all such points $P$, the quantity $PA+PB$ attains its minimum value at a unique point $P^{*}$. The value of $PA\cdot PB$ at $P=P^{*}$ is:
\begin{enumerate}[label=(\Alph*),leftmargin=*]
  \item $36$
  \item $40$
  \item $41$
  \item $50$
\end{enumerate}

\vspace{0.6em}\hrule\vspace{0.6em}

%================================================================
\section*{Answer Key}

\begin{center}
\renewcommand{\arraystretch}{1.25}
\begin{tabular}{|c|l|c|c|}
\hline
\textbf{Q.No.} & \textbf{Topic} & \textbf{Type} & \textbf{Answer} \\
\hline
1  & Arithmetic -- TSD                     & MCQ  & (B) $18.5$ \\
2  & Arithmetic -- Mixtures                & TITA & $20$ \\
3  & Arithmetic -- Time and Work           & MCQ  & (B) $36$ \\
4  & Arithmetic -- Profit and Loss         & TITA & $20$ \\
5  & Arithmetic -- Averages                & MCQ  & (B) $3:2$ \\
6  & Algebra -- Functional Equations       & TITA & $-32$ \\
7  & Algebra -- Logarithms                 & MCQ  & (C) $2$ \\
8  & Algebra -- Inequalities               & MCQ  & (A) $4$ \\
9  & Algebra -- Quadratics                 & MCQ  & (B) $2$ \\
10 & Number System -- CRT                  & TITA & $1001$ \\
11 & Number System -- Divisibility         & MCQ  & (C) $51$ \\
12 & Number System -- Digits               & TITA & $825$ \\
13 & Geometry -- Cevians                   & MCQ  & (A) $3:5$ \\
14 & Mensuration -- Tangents               & TITA & $30$ \\
15 & Geometry -- Optimisation              & MCQ  & (C) $41$ \\
\hline
\end{tabular}
\end{center}

\vspace{0.4em}\hrule\vspace{0.6em}

%================================================================
\section*{Detailed Solutions}

\subsection*{Solution 1 -- Two Cyclists, Two Meetings}
Let $D$ be the distance $AB$ and let the speeds of $P, Q$ be $v_P, v_Q$.

At the first meeting, $P$ has covered $7.5$ and $Q$ has covered $D-7.5$, both in the same time. So
\[
\frac{v_P}{v_Q}=\frac{7.5}{D-7.5}. \tag{1}
\]

By the time of the second meeting, both cyclists have reached the far town and turned back (this must be checked at the end). Between start and second meeting, the two together cover $3D$ (the first meeting accounts for a combined $D$; after that, until they meet again, they together cover an additional $2D$).

In total, $P$ travels $D$ to reach $B$, then turns and is at $4$ km from $B$, so $P$'s total distance is $D + 4$. Hence $Q$'s total distance is $3D - (D+4) = 2D - 4$.

\[
\frac{v_P}{v_Q}=\frac{D+4}{2D-4}. \tag{2}
\]

Equating (1) and (2):
\[
\frac{7.5}{D-7.5}=\frac{D+4}{2D-4}
\;\Longrightarrow\; 7.5(2D-4)=(D+4)(D-7.5).
\]
$15D-30 = D^2 - 3.5D - 30 \;\Longrightarrow\; D^2 - 18.5D = 0 \;\Longrightarrow\; D = 18.5.$

\textbf{Verification of the ``both have turned'' assumption:} With $D=18.5$, speed ratio $v_P:v_Q = 7.5:11 = 15:22$. Time for $P$ to reach $B$ is $18.5/15 \approx 1.233$; time for $Q$ to reach $A$ is $18.5/22 \approx 0.841$. Time of second meeting solves $37 - 15t = 22t - 18.5$, giving $t=1.5$, which exceeds both turnaround times. Position $= 14.5$, i.e., $4$ from $B$. \checkmark

\textbf{Why this is CAT-level:} The natural first approach (``ratio of distances at first meeting and total covered $=3D$'') is correct only if both cyclists have turned by the second meeting -- a hidden case-check that distinguishes a percentile-pushing solver.

\textbf{Answer: (B) $18.5$.}

\medskip

\subsection*{Solution 2 -- Mixtures with Integer Replacement}
After $4$ operations, the milk remaining is
\[
80\!\left(\frac{80-x}{80}\right)^{4}=\frac{(80-x)^{4}}{80^{3}}.
\]
We need $25<\dfrac{(80-x)^{4}}{80^{3}}<30$, i.e., $25\cdot 80^{3}<(80-x)^{4}<30\cdot 80^{3}.$

$80^3 = 512\,000$, so $1.28\times 10^{7}<(80-x)^{4}<1.536\times 10^{7}.$

Test integers $80-x$:
\begin{itemize}[leftmargin=*,itemsep=0pt]
  \item $80-x=59:\ 59^{4}=12{,}117{,}361$ -- too small.
  \item $80-x=60:\ 60^{4}=12{,}960{,}000$ -- in range ($x=20$, milk $\approx 25.31$).
  \item $80-x=61:\ 61^{4}=13{,}845{,}841$ -- in range ($x=19$, milk $\approx 27.04$).
  \item $80-x=62:\ 62^{4}=14{,}776{,}336$ -- in range ($x=18$, milk $\approx 28.86$).
  \item $80-x=63:\ 63^{4}=15{,}752{,}961$ -- too large.
\end{itemize}

So $x\in\{18,19,20\}$. The \textbf{largest} such $x$ is $\boxed{20}$.

\textbf{Why this is CAT-level:} The hardening comes from the band condition combined with integer feasibility. The natural attempt -- setting up $(80-x)^4 = k\cdot 80^3$ for a clean $k$ -- fails because no clean closed-form root works; one must \emph{bracket} the integer band by estimating $60^4$, $61^4$, $62^4$ mentally.

\textbf{Answer: $x = 20$.}

\medskip

\subsection*{Solution 3 -- Time and Work}
Combined rate of A and B $=1/20$. B's rate $=1/45$.

So A's rate $=1/20 - 1/45 = (9-4)/180 = 5/180 = 1/36.$

Therefore A alone takes $\mathbf{36}$ days.

\textbf{Consistency check with the third datum:} Let A and B work together for $t$ days, then B alone for $21.6$ days. Total work: $t/20 + 21.6/45 = 1 \Rightarrow t/20 = 1 - 0.48 = 0.52 \Rightarrow t = 10.4$ days. Consistent.

\textbf{Why this is CAT-level:} The redundant ``$21.6$ days'' datum tempts students into setting up a two-variable equation system. The elegant route uses only the first two facts, with the third merely confirming consistency. Method selection -- not method execution -- is what saves time.

\textbf{Answer: (B) $36$.}

\medskip

\subsection*{Solution 4 -- Faulty Balance with Discount}
Let CP $= \text{Rs.}\,100$ per $1000$ g. Then MP $= \text{Rs.}\,140$ per displayed $1000$ g. SP (after discount $d\%$) $= 140(1-d/100)$ per displayed $1000$ g.

The trader actually delivers $800$ g per displayed $1000$ g. The CP of $800$ g $=\text{Rs.}\,80$. The trader receives $140(1-d/100)$.

Given: profit $= 40\%$ of actual CP, i.e., profit $= 0.40\times 80 = 32$. So
\[
140(1-d/100) - 80 = 32 \;\Longrightarrow\; 140(1-d/100) = 112 \;\Longrightarrow\; 1-d/100 = 0.8 \;\Longrightarrow\; d = 20.
\]

\textbf{Why this is CAT-level:} The trap is the ambiguity in ``base of profit.'' Profit measured against \emph{displayed-weight CP} vs.\ \emph{actual delivered CP} gives different answers. The clean reading -- ``profit on the goods physically delivered'' -- fixes the base. Students who default to displayed-weight CP get a different (wrong) answer.

\textbf{Answer: $d = 20$.}

\medskip

\subsection*{Solution 5 -- Averages and Ratios}
Let $B$, $G$ be the original numbers of boys and girls. Original average:
\[
\frac{68B + 80G}{B+G} = 74 \;\Longrightarrow\; 68B + 80G = 74B + 74G \;\Longrightarrow\; 6G = 6B \;\Longrightarrow\; B=G.
\]

Let $n$ new boys (average $62$) join. New overall average:
\[
\frac{74(B+G) + 62n}{B+G+n} = 71 \;\Longrightarrow\; 74\cdot 2B + 62n = 71(2B+n) \;\Longrightarrow\; 148B + 62n = 142B + 71n.
\]
$6B = 9n \;\Longrightarrow\; B/n = 9/6 = 3/2.$

\textbf{Original boys : new boys $= 3:2$.}

\textbf{Why this is CAT-level:} The two-stage weighted-average computation forces students to derive $B=G$ as an interim result, then use it. A direct alligation attempt on the new group vs the original mean (gap $74-62=12$, gap $74-71=3$) gives ratio of (original total) : (new) $= 12:3 = 4:1$, i.e., $2B:n = 4:1 \Rightarrow B = 2n$ -- a tempting trap that ignores the original boy-girl split. Careful weighted averaging recovers $B = 1.5n$, i.e., $3:2$.

\textbf{Answer: (B) $3:2$.}

\medskip

\subsection*{Solution 6 -- Functional Equation}
Given $f(x)+2\,f(1-x)=x^{2}+3x.$ Replace $x \to 1-x$:
\[
f(1-x)+2\,f(x)=(1-x)^{2}+3(1-x)=x^{2}-5x+4.
\]
Let $u=f(x),\ v=f(1-x)$. System: $u+2v=x^{2}+3x$, $2u+v=x^{2}-5x+4.$

$2\cdot\text{(second)} - \text{(first)}: 3u = 2x^{2}-10x+8 - (x^{2}+3x) = x^{2}-13x+8.$

So $f(x) = \dfrac{x^{2}-13x+8}{3}.$

$f(5) = \dfrac{25-65+8}{3} = \dfrac{-32}{3}.$ Hence $3\,f(5) = -32.$

\textbf{Verification:} $f(5) + 2f(-4) = -32/3 + 2\cdot (16+52+8)/3 = -32/3 + 152/3 = 120/3 = 40 = 25 + 15.$ \checkmark

\textbf{Why this is CAT-level:} The natural false start -- guessing $f$ is a quadratic and matching coefficients -- works but is slower. The clean route uses the substitution $x\to 1-x$, but students must recognise that the involution $x \mapsto 1-x$ is the right swap (not $x \to -x$). Method selection under time pressure is the test.

\textbf{Answer: $3\,f(5) = -32.$}

\medskip

\subsection*{Solution 7 -- Logarithmic Equation}
LHS $=\log_{2}[(x^{2}-3x+2)(x^{2}-7x+12)]= \log_2[(x-1)(x-2)(x-3)(x-4)].$

RHS $=\log_2 8 + \log_2\!\left[\dfrac{(x-1)(x-4)}{2}\right] = \log_2\!\left[\dfrac{8(x-1)(x-4)}{2}\right] = \log_2[4(x-1)(x-4)].$

\textbf{Domain:}
$(x-1)(x-2)>0,\ (x-3)(x-4)>0,\ (x-1)(x-4)>0.$
Intersection: $x<1$ or $x>4.$

On this domain, $(x-1)(x-4)>0$, so the equation becomes
\[
(x-1)(x-2)(x-3)(x-4) = 4(x-1)(x-4) \;\Longrightarrow\; (x-2)(x-3) = 4.
\]
$x^{2}-5x+6=4 \;\Longrightarrow\; x^{2}-5x+2=0 \;\Longrightarrow\; x = \dfrac{5\pm\sqrt{17}}{2}.$

$\dfrac{5+\sqrt{17}}{2}\approx 4.56 > 4$ \checkmark and $\dfrac{5-\sqrt{17}}{2}\approx 0.44 < 1$ \checkmark. Both are in the domain.

So there are \textbf{2 real values}.

\textbf{Why this is CAT-level:} The temptation is to immediately set $(x-2)(x-3)=4$ without checking the log domain. Students who skip the domain analysis still arrive at $2$ -- but by accident. A natural error path is to assume one of the roots is rejected by the domain; only careful interval analysis confirms both lie in the valid region.

\textbf{Answer: (C) $2$.}

\medskip

\subsection*{Solution 8 -- Maximum of $a$ under Two Constraints}
Given $a+b+c=12$ and $abc=64$ with $b,c>0$. For fixed $a\in(0,12)$, $b+c = 12-a$ and $bc = 64/a$. For $b,c$ to be real and positive, we need
\[
(b+c)^{2}\ge 4bc \;\Longrightarrow\; (12-a)^{2}\ge \frac{256}{a} \;\Longrightarrow\; a(12-a)^{2}\ge 256. \tag{*}
\]
Define $g(a)=a(12-a)^{2}$ on $(0,12)$. Then
\[
g'(a)=(12-a)^{2}+a\cdot 2(12-a)(-1)=(12-a)[(12-a)-2a]=(12-a)(12-3a).
\]
$g'(a)=0$ at $a=12$ (boundary) and $a=4$. $g$ is increasing on $(0,4)$ and decreasing on $(4,12)$. Maximum at $a=4$: $g(4)=4\cdot 64 = 256$.

So $g(a)\le 256$ on $(0,12)$ with equality \emph{only} at $a=4$. Thus the inequality (*) is satisfied only at $a=4$.

At $a=4$: $b+c=8, bc=16 \Rightarrow b=c=4.$

\textbf{Maximum possible $a = 4.$}

\textbf{Why this is CAT-level:} The natural approach -- AM-GM gives $a+b+c\ge 3\sqrt[3]{abc}$, i.e., $12\ge 3\sqrt[3]{64}=12$ -- forces equality, so $a=b=c=4$. This collapses the system to a single point. The trap is to instead bound $a$ from above by trying to make $b,c$ very small, missing that equality in AM-GM is forced by the data. The discriminant route (*) confirms uniqueness rigorously. Option (B) and (D) reflect false bounds derived by trying $b=c$ but solving an incorrect equation.

\textbf{Answer: (A) $4$.}

\medskip

\subsection*{Solution 9 -- Quadratic with Linked Condition}
By Vieta's: $p+q=k+3,\ pq=2k+1.$ 
\[
(p-q)^{2}=(p+q)^{2}-4pq=(k+3)^{2}-4(2k+1)=k^{2}+6k+9-8k-4=k^{2}-2k+5.
\]
Set $(p-q)^{2}=12$:
\[
k^{2}-2k+5=12 \;\Longrightarrow\; k^{2}-2k-7=0.
\]
Discriminant $=4+28=32>0$, real distinct roots. Sum of roots $=2.$

\textbf{Verification of root reality:} For each such $k$, $(p-q)^2 = 12 > 0$, so $p, q$ are real and distinct. Both values of $k$ are admissible.

\textbf{Why this is CAT-level:} Mild but real. The standard Vieta route is direct, but the natural error is forgetting the $4pq$ term sign or confusing $(p-q)^2$ with $(p+q)^2$. The wrong options trap exactly those errors: option (A) $-2$ corresponds to summing roots of $k^2+2k-7$ (sign slip on the linear term).

\textbf{Answer: (B) $2$.}

\medskip

\subsection*{Solution 10 -- CRT with Filtering}
$N\equiv 3\pmod 7,\ N\equiv 5\pmod{11},\ N\equiv 0\pmod{13}.$

Since $7,11,13$ are pairwise coprime, by the Chinese Remainder Theorem, $N$ is uniquely determined modulo $7\cdot 11\cdot 13 = 1001.$

\textbf{Finding $N$:} Write $N=13k$. Then
\begin{align*}
13k \equiv 3\pmod 7 &\Rightarrow 6k \equiv 3 \Rightarrow -k \equiv 3 \Rightarrow k \equiv 4\pmod 7,\\
13k \equiv 5\pmod{11} &\Rightarrow 2k \equiv 5 \Rightarrow k \equiv 5\cdot 6 = 30 \equiv 8 \pmod{11}.
\end{align*}
($2^{-1}\equiv 6 \pmod{11}$ since $2\cdot 6 = 12\equiv 1$.)

Combine: $k=7m+4$, then $7m+4\equiv 8\pmod{11} \Rightarrow 7m\equiv 4 \Rightarrow m\equiv 4\cdot 8=32\equiv 10\pmod{11}$ (since $7\cdot 8 = 56\equiv 1\pmod{11}$).

So $m = 11j+10$, $k = 77j+74$. Smallest positive $k$ is $74$, giving $N = 13\cdot 74 = 962.$

\textbf{Verification:} $962 = 7\cdot 137 + 3$ \checkmark, $962 = 11\cdot 87 + 5$ \checkmark, $962/13 = 74$ \checkmark.

The next solution is $N + 1001 = 1963$. So $M = 1963$ and $M-N = 1001.$

\textbf{Why this is CAT-level:} The question \emph{looks} like a standard CRT problem, but the third condition (``divisible by $13$'') is the moduli twist -- it must be expressed as $N\equiv 0\pmod{13}$. Two layers of modular inversion ($2^{-1}\pmod{11}$, $7^{-1}\pmod{11}$) are required. The clean route exploits the fact that $M-N$ is simply $\mathrm{lcm}(7,11,13)=1001$, sidestepping the entire computation of $N$ itself. Students who don't see this shortcut compute $N$ first and then add the LCM -- which works but is slower.

\textbf{Answer: $M - N = 1001.$}

\medskip

\subsection*{Solution 11 -- Exactly Two of $\{4,6,9\}$}
Let $A, B, C$ denote the sets of multiples of $4, 6, 9$ in $\{1,2,\ldots,600\}$.
\begin{itemize}[leftmargin=*,itemsep=0pt]
  \item $\mathrm{lcm}(4,6)=12,\ \mathrm{lcm}(4,9)=36,\ \mathrm{lcm}(6,9)=18,\ \mathrm{lcm}(4,6,9)=36.$
  \item $|A\cap B|=\lfloor 600/12\rfloor=50,\ |A\cap C|=\lfloor 600/36\rfloor=16,\ |B\cap C|=\lfloor 600/18\rfloor=33,\ |A\cap B\cap C|=16.$
\end{itemize}

Count divisible by exactly two: $\sum_{\text{pairs}}|A_i\cap A_j| - 3|A\cap B\cap C|$
\[
=(50+16+33)-3\cdot 16 = 99 - 48 = 51.
\]

\textbf{Structural insight:} Since $\mathrm{lcm}(4,9)=36=\mathrm{lcm}(4,6,9)$, every multiple of both $4$ and $9$ is automatically a multiple of $6$. So the pair $\{4,9\}$ alone is impossible -- contributing $0$. The actual count breaks as $(50-16) + 0 + (33-16) = 34 + 17 = 51.$

\textbf{Why this is CAT-level:} The trap is mechanical inclusion--exclusion without noticing the structural collapse $\mathrm{lcm}(4,9)=\mathrm{lcm}(4,6,9)$. Students who plug into a memorised formula get the right number ($51$) accidentally; students who reason from structure get it confidently. Wrong option (B) $34$ is the count from the $\{4,6\}$ pair alone -- a natural partial answer.

\textbf{Answer: (C) $51$.}

\medskip

\subsection*{Solution 12 -- Three-Digit with Three Linked Conditions}
Write $\overline{abc}=100a+10b+c$ with $1\le a\le 9$, $0\le b\le 9$, $1\le c\le 9$ (the last because $\overline{cba}$ is a three-digit number).

\textbf{Condition (iii):} $\overline{cba}>\overline{abc} \Leftrightarrow 99(c-a)>0 \Leftrightarrow c>a.$

\textbf{Condition (i)} (divisibility by $11$): $a - b + c \equiv 0\pmod{11}.$ Since $a,b,c\in[0,9]$, $a-b+c\in[-9,18]$. The achievable multiples of $11$ are $0$ and $11.$

\textbf{Case 1: $a-b+c=0$, i.e., $b=a+c.$}

Combine with condition (ii): $a+b+c=14 \Rightarrow 2(a+c)=14 \Rightarrow a+c=7,\ b=7.$

With $c>a$, $a\ge 1$, $c\le 9$: pairs $(a,c)$ are $(1,6),(2,5),(3,4)$.

Numbers: $176,\ 275,\ 374.$

\textbf{Case 2: $a-b+c=11$, i.e., $b = a+c - 11.$}

Combine: $a+b+c = 2(a+c)-11 = 14 \Rightarrow a+c = 12.5$, not an integer. No solutions.

\textbf{Sum} $= 176+275+374 = 825.$

\textbf{Why this is CAT-level:} Three conditions interact non-trivially. Condition (iii) -- expressed as ``$\overline{cba}-\overline{abc}$ is a positive multiple of $99$'' -- looks restrictive but is automatic for any reversal; the actual content is just $c>a$. Recognising this collapse is the key insight. The divisibility rule for $11$ produces a case split (Case 2 is killed by parity).

\textbf{Answer: $176+275+374 = 825.$}

\medskip

\subsection*{Solution 13 -- Cevians and Mass Points}
Place $A$ at the origin and use vectors. Let $\vec b = \vec{AB}, \vec c = \vec{AC}.$

$D$ on $BC$ with $BD:DC=2:3$: $\vec{AD} = \dfrac{3\vec b + 2\vec c}{5}.$

$E$ on $AD$ with $AE:ED=3:2$: $\vec{AE} = \dfrac{3}{5}\vec{AD} = \dfrac{3(3\vec b + 2\vec c)}{25} = \dfrac{9\vec b + 6\vec c}{25}.$

Parametrise line $BE$: $\vec r = \vec b + t(\vec{AE} - \vec b) = \vec b + t\!\left(\dfrac{9\vec b + 6\vec c}{25} - \vec b\right) = \vec b\!\left(1 - \dfrac{16t}{25}\right) + \vec c\cdot\dfrac{6t}{25}.$

For this point to lie on line $AC$ (i.e., $\vec r = s\vec c$), the coefficient of $\vec b$ must vanish:
\[
1 - \dfrac{16t}{25} = 0 \;\Longrightarrow\; t = \dfrac{25}{16}.
\]
Then $\vec r = \vec c\cdot\dfrac{6}{25}\cdot\dfrac{25}{16} = \dfrac{3}{8}\vec c.$

So $\vec{AF} = \dfrac{3}{8}\vec c$, meaning $AF:AC = 3:8$, hence
\[
AF:FC = 3:5.
\]

\textbf{Cross-check by mass points:} Place masses: at $C$ mass $2$, at $B$ mass $3$ so $D$ has mass $5$ (since $BD:DC=2:3$). Place mass at $A$ such that $AE:ED=3:2$, i.e., mass at $A$ should be in ratio $2:3$ with mass at $D$, so mass at $A$ is $2$ and mass at $D$ is $3$ (consistent up to scale; rescale $D$'s mass from $5$ to $3$ by dividing by $5/3$; equivalently, multiply $A$'s mass by $5/3$). To keep arithmetic clean: mass at $A=2\cdot 5 = 10$, at $D = 3\cdot 5 = 15$, so at $B$ effective mass $= (3/5)\cdot 15 = 9$ and at $C$ effective mass $= (2/5)\cdot 15 = 6.$ Then on $AC$, masses $10$ at $A$ and $6$ at $C$, so $F$ divides $AC$ in ratio $6:10 = 3:5$ from $A$. \checkmark

\textbf{Why this is CAT-level:} The natural temptation is to use Menelaus on $\triangle ADC$ with transversal $B$--$E$--$F$, but choosing the right triangle and getting the signed ratios correct is error-prone under time pressure. Vector / mass-point methods are faster and trap-resistant. Wrong option (C) $9:10$ corresponds to a common Menelaus sign error.

\textbf{Answer: (A) $3:5$.}

\medskip

\subsection*{Solution 14 -- Common External Tangent}
Two circles of radii $R=5$ and $r=3$, externally tangent (so the distance between centres is $R+r=8$). The length of a common external tangent between two circles with centre-distance $d$ and radii $R, r$ is
\[
L=\sqrt{d^{2}-(R-r)^{2}}.
\]
Here $L = \sqrt{8^{2} - 2^{2}} = \sqrt{64 - 4} = \sqrt{60} = 2\sqrt{15}.$

\textbf{Encoded answer:} $a=2,\ b=15,\ a\cdot b = 30.$

\textbf{Why this is CAT-level:} The standard formula is two distinct objects: external tangent uses $(R-r)$, internal tangent uses $(R+r)$. Many students confuse the two -- especially when the circles are already externally tangent (which makes the internal tangent degenerate). Choosing the correct formula \emph{is} the test. The encoded answer format ($a\cdot b = 30$) blocks accidental matches with $\sqrt{60}$, $4\sqrt{15}$, etc.

\textbf{Answer: $a\cdot b = 30.$}

\medskip

\subsection*{Solution 15 -- Minimising $PA+PB$ on a Locus}
Area constraint: $[\triangle APB] = \tfrac{1}{2}\cdot AB \cdot y = \tfrac{1}{2}\cdot 10\cdot y = 5y = 20 \Rightarrow y = 4.$

So $P$ lies on the horizontal line $y=4$, with $P=(x,4)$ for some real $x$.

\textbf{Minimising $PA+PB$ on the line $y=4$:} Reflect $A=(0,0)$ across this line to $A'=(0,8)$. For any $P$ on the line, $PA = PA'$, so
\[
PA + PB = PA' + PB \ge A'B.
\]
Equality iff $A', P, B$ are collinear. $A'B = \sqrt{(10-0)^{2} + (0-8)^{2}} = \sqrt{164} = 2\sqrt{41}.$

Line $A'B$: from $(0,8)$ to $(10,0)$, slope $-0.8$, equation $y = 8 - 0.8x.$ At $y=4$: $x = 5.$ So $P^{*} = (5,4).$

\textbf{Computing $PA\cdot PB$ at $P^{*}$:} By symmetry of $(5,4)$ relative to $A=(0,0)$ and $B=(10,0)$:
\[
PA = \sqrt{25+16} = \sqrt{41}, \qquad PB = \sqrt{25+16} = \sqrt{41}.
\]
Hence $PA\cdot PB = 41.$

\textbf{Why this is CAT-level:} The reflection trick is the canonical fast route, but the question disguises it -- students may try Lagrange multipliers, derivatives, or even ellipse properties (since $PA+PB$ constant defines an ellipse with foci $A,B$, and we want the smallest ellipse touching $y=4$). All routes converge, but the reflection method is by far the cleanest under time pressure. Wrong option (D) $50$ tempts those who compute $PA \cdot PB$ at $P=(0,4)$ (a boundary candidate, $\sqrt{16}\cdot\sqrt{116}\approx 4\cdot 10.77 = 43$, close to but not $50$); wrong option (A) $36$ corresponds to $P=(5,4)$ but with an area constraint of $25$ (a misreading).

\textbf{Answer: (C) $41$.}

\vspace{0.6em}\hrule\vspace{0.6em}

%================================================================
\section*{Topic-Wise Distribution}

\begin{center}
\renewcommand{\arraystretch}{1.25}
\begin{tabular}{|l|c|c|c|c|c|}
\hline
\textbf{Topic} & \textbf{Count} & \textbf{Medium-Hard} & \textbf{Hard} & \textbf{Very Hard} & \textbf{MCQ : TITA} \\
\hline
Arithmetic           & 5 & 1 & 3 & 1 & 3 : 2 \\
Algebra              & 4 & 1 & 2 & 1 & 3 : 1 \\
Number System        & 3 & 0 & 2 & 1 & 1 : 2 \\
Geometry/Mensuration & 3 & 1 & 1 & 1 & 2 : 1 \\
\hline
\textbf{Total}       & \textbf{15} & \textbf{3} & \textbf{8} & \textbf{4} & \textbf{9 : 6} \\
\hline
\end{tabular}
\end{center}

\textit{Note: The distribution shows 3 Medium-Hard, 8 Hard, and 4 Very Hard. Q14 is Hard but borderline Very-Hard due to the tangent-type discrimination trap; counted here under Hard.}

\vspace{0.4em}\hrule\vspace{0.6em}

%================================================================
\section*{Quality Verification Summary}

\begin{center}
\renewcommand{\arraystretch}{1.2}
\small
\begin{tabular}{|c|c|c|c|c|c|c|}
\hline
\textbf{Q} & \textbf{Answer Verified} & \textbf{MCQ Options Verified} & \textbf{Ambiguity Check} & \textbf{Brute Force Resistance} & \textbf{Method Selection Needed} & \textbf{Status} \\
\hline
1  & Yes & Yes            & Clear wording & Yes      & Yes & Ready \\
2  & Yes & Not Applicable & Clear wording & Yes      & Yes & Ready \\
3  & Yes & Yes            & Clear wording & Yes      & Yes & Ready \\
4  & Yes & Not Applicable & Clear wording & Moderate & Yes & Ready \\
5  & Yes & Yes            & Clear wording & Yes      & Yes & Ready \\
6  & Yes & Not Applicable & Clear wording & Yes      & Yes & Ready \\
7  & Yes & Yes            & Clear wording & Yes      & Yes & Ready \\
8  & Yes & Yes            & Clear wording & Yes      & Yes & Ready \\
9  & Yes & Yes            & Clear wording & Moderate & Yes & Ready \\
10 & Yes & Not Applicable & Clear wording & Yes      & Yes & Ready \\
11 & Yes & Yes            & Clear wording & Yes      & Yes & Ready \\
12 & Yes & Not Applicable & Clear wording & Yes      & Yes & Ready \\
13 & Yes & Yes            & Clear wording & Yes      & Yes & Ready \\
14 & Yes & Not Applicable & Clear wording & Yes      & Yes & Ready \\
15 & Yes & Yes            & Clear wording & Yes      & Yes & Ready \\
\hline
\end{tabular}
\end{center}

\end{document}
